\chapter{Resumen de la Semana}
Este documento consolida los avances tecnicos en el desarrollo, calibracion y validacion metrologica de un sistema de microscopia computacional basado en sensor CMOS, concebido como carga util experimental (payload) para el proyecto INTISAT.

El objetivo principal del payload es demostrar la viabilidad de un sistema de caracterizacion de muestras biologicas y biomimeticas mediante una arquitectura compacta, modular y de bajo consumo, adecuada para plataformas con recursos limitados. Para ello se emplea una estrategia de adquisicion de imagenes basada en hardware minimo y procesamiento computacional.

Como etapa central del trabajo, se aborda la validacion de la escala espacial ($\mu$m/pixel) del sistema de imagen. Se utilizan microesferas de poliestireno de diametro nominal de 2 $\mu$m como objetos patron, comparando mediciones del sistema CMOS con referencias de microscopia optica y microscopia electronica de barrido (SEM).

La validacion busca asegurar coherencia entre dimensiones fisicas reales y mediciones digitales, considerando limites de resolucion del sensor, optica de contacto y efectos de difraccion. Este paso es previo a una caracterizacion mas robusta de muestras biologicas reales.

Finalmente, se documentan pruebas, resultados preliminares y lineas de trabajo para las siguientes etapas de integracion y mejora del payload.